% ═══════════════════════════════════════════════════════════════════
% ПРИЛОЖЕНИЯ K–L: ОРБИТЫ ПРЯМЫХ И СПРАВОЧНИК КОМАНД
% ═══════════════════════════════════════════════════════════════════

\section{Орбитная структура $48$ прямых K3}\label{app:orbits}

Действие $(\mu_4)^2$ на прямых имеет $12$ орбит по $4$ прямых.
Ниже --- структура орбит и вклад в кратности характеров
(сертификат~C, теорема~\ref{th:C1}).

\subsection{Семейства и орбиты}

Каждое семейство $L_f$ ($f=1,2,3$) содержит $16$ прямых ---
сетку $(a,b)\in(\ZZ/4)^2$. Действие поворотов сдвигает индексы:
орбиты --- циклы по одному из индексов при фиксированном другом.
Полная структура:

\begin{center}
\small
\begin{tabular}{@{}l l l l@{}}
\toprule
Семейство & Орбита & Прямые & Характер \\
\midrule
$L_1$ & $\{(a,b): b=a\}$-циклы & $4$ орбиты по $4$ &
  смешанные \\
$L_1$ & сдвиги $a$ при $b$ фикс. & $4$ орбиты по $4$ &
  смешанные \\
$L_2$ & аналогично & $4+4$ & смешанные \\
$L_3$ & аналогично & $4+4$ & смешанные \\
\bottomrule
\end{tabular}
\end{center}

Проекционные операторы на характеры дают точные кратности:
тривиальный характер --- $1$ (диагональная комбинация сумм
семейств), каждый нетривиальный --- $7$. Проверка через вторую
схему: спектр перестановочной матрицы $48\times48$ имеет
собственные значения $1$ (кратность $1$), $7i$, $-7$, $-7i$ как
суммарные следы блоков --- согласовано.

\subsection{Согласование с инерцией}

Двенадцать орбит по четыре прямых дают $48$; добавление $h$
замыкает $49$ измерений; ранг $20$, инерция $(1,19)$ ---
соответствие таблицы орбит с числовой сигнатурой --- ещё одна
двойная регистрация (комбинаторика $\leftrightarrow$ линейная
алгебра).

\section{Справочник команд лаборатории}\label{app:cli}

\subsection{Запуск}

\begin{center}
\small
\begin{tabularx}{\textwidth}{@{}l >{\raggedright\arraybackslash}X@{}}
\toprule
Команда & Действие \\
\midrule
\texttt{python3 laboratory.py} & интерактивное меню (выбор языка,
полный доступ) \\
\texttt{python3 laboratory.py --lang en} & старт с английским
интерфейсом \\
\texttt{python3 laboratory.py --lang ru} & старт с русским
интерфейсом \\
\texttt{python3 laboratory.py --run all} & неинтерактивный полный
протокол \\
\texttt{python3 laboratory.py --run v1-v9} & только протокол
V1--V9 \\
\texttt{python3 laboratory.py --run stands} & все стенды \\
\texttt{python3 laboratory.py --run certs} & все сертификаты \\
\texttt{python3 laboratory.py --run plots} & только диаграммы \\
\texttt{python3 laboratory.py --no-plots} & прогон без графики
(минимальные зависимости) \\
\texttt{python3 laboratory.py --dps 50} & повышенная точность \\
\bottomrule
\end{tabularx}
\end{center}

\subsection{Структура меню}

\begin{center}
\small
\begin{tabularx}{\textwidth}{@{}l >{\raggedright\arraybackslash}X@{}}
\toprule
Пункт & Содержание \\
\midrule
1 & Полный протокол V1--V9 \\
2 & Стенды: тор / K3 / Клейн / errata E8 / бинарный код \\
3 & Сертификаты A--H (по выбору или все) \\
4 & Стенд N=15/30: замкнутая система (V1--V9 + периоды) \\
5 & Параметры расчётов: $dps$, выборки, уровни $N$, число тестов \\
6 & Конструктор экспериментов: произвольные $N$, $(a,b)$,
  $(r,s)$; переписи; CM-решётки; время завершимости \\
7 & Диаграммы: плиточная система $600$~dpi (восемь плиток) \\
8 & Отчёты: \texttt{reports/} --- JSON + журнал + плитки \\
9 & Мультиязычная верификация: Julia / Fortran / C / Rust / Lean \\
0 & О программе, лицензия, цитирование \\
L & Переключение языка (RU/EN) в любой момент \\
\bottomrule
\end{tabularx}
\end{center}

\subsection{Параметры конструктора экспериментов}

\begin{center}
\small
\begin{tabularx}{\textwidth}{@{}l l >{\raggedright\arraybackslash}X@{}}
\toprule
Параметр & Диапазон & Смысл \\
\midrule
$N$ & $4\dots64$ & уровень кривой Ферма \\
$(a,b)$ & $1\le a,b$, $a+b\le N-1$ & характер \\
$(r,s)$ & $0\dots N-1$ & данные обхода \\
$dps$ & $20\dots120$ & точность \texttt{mpmath} \\
$W,H$ & $\le 512$ & сетка потока (E) \\
$d$ & $1\dots100$ & CM-дискриминант решётки (G) \\
тестов/проводник & $1\dots5$ & плотность выборки V2 \\
\bottomrule
\end{tabularx}
\end{center}

Конструктор печатает замкнутую форму, независимое интегрирование,
относительное отклонение и вердикт --- полный мини-сертификат
для произвольной конфигурации; результаты добавляются в отчёт
запуска.

\subsection{Коды возврата и сверка}

Код возврата $0$ --- все проверки PASS; $1$ --- есть FAIL;
$2$ --- отсутствует зависимость. CI (GitHub Actions) использует
коды возврата: workflow \texttt{ci.yml} считает сборку зелёной
только при \texttt{--run all} с кодом $0$ на трёх версиях Python.
